% from ega4-16-fr.tex
\oldpage[IV-4]{33}
affines $V'\subset V$ et $W'\subset W$. Mais en vertu de la description
précédente de $h$, cela résulte de la commutativité du diagramme (0, 20.5.4.2).

Les applications
$\Gamma(U_0,\mathscr{G})\times P(U_0)\to P(U_0)$ définissent donc un
homomorphisme de \emph{faisceaux d'ensembles}
$m:\mathscr{G}\times\mathscr{P}\to\mathscr{P}$ tel que, pour tout ouvert
$U_0$ pour lequel $\Gamma(U_0,\mathscr{P})\neq\varnothing$,
\[
  m_{U_0}:\Gamma(U_0,\mathscr{G})\times\Gamma(U_0,\mathscr{P})
  \longrightarrow\Gamma(U_0,\mathscr{P})
\]
soit une loi externe définissant sur $\Gamma(U_0,\mathscr{P})$ une structure
de torseur sous le groupe $\Gamma(U_0,\mathscr{G})$.
\end{env}

\begin{env}[16.5.15]
\EGAReaderTarget{ega4.c04.statement.16-5-15}{IV.16.5.15}
D'une façon générale, lorsqu'on se donne sur un espace topologique $Z$ un
faisceau d'ensembles $\mathscr{P}$, un faisceau de groupes $\mathscr{G}$ (non
nécessairement commutatif) et un homomorphisme de faisceaux d'ensembles
$m:\mathscr{G}\times\mathscr{P}\to\mathscr{P}$ tel que, pour tout ouvert
$U\subset Z$ tel que $\Gamma(U,\mathscr{P})\neq\varnothing$,
\[
  m_U:\Gamma(U,\mathscr{G})\times\Gamma(U,\mathscr{P})
  \longrightarrow\Gamma(U,\mathscr{P})
\]
fasse de $\Gamma(U,\mathscr{P})$ un \emph{torseur} sous le groupe
$\Gamma(U,\mathscr{G})$, on dit que $\mathscr{P}$ est un
\emph{pseudo-torseur} (ou \emph{faisceau formellement principal homogène})
sous le faisceau de groupes $\mathscr{G}$. On dit que $\mathscr{P}$ est un
\emph{torseur} (ou \emph{faisceau principal homogène}) sous $\mathscr{G}$ si
de plus $\Gamma(U,\mathscr{P})\neq\varnothing$ pour tout ouvert
$U\neq\varnothing$ d'une base convenable de la topologie de $Z$.

Pour la théorie générale des torseurs, nous renvoyons à [42]; nous nous
bornerons ici à rappeler la correspondance canonique entre les classes
d'isomorphie de ces torseurs (pour un $\mathscr{G}$ donné) et les éléments de
l'ensemble de cohomologie $H^1(Z,\mathscr{G})$. Considérons en effet un
torseur $\mathscr{P}$ sous $\mathscr{G}$ et un recouvrement ouvert
$(U_\lambda)$ de $Z$ tel que
$\Gamma(U_\lambda,\mathscr{P})\neq\varnothing$ pour tout $\lambda$;
désignons par $p_\lambda$ un élément de
$\Gamma(U_\lambda,\mathscr{P})$. Pour tout couple d'indices $\lambda$, $\mu$
tel que $U_\lambda\cap U_\mu\neq\varnothing$, il existe alors un élément et
un seul $\gamma_{\lambda\mu}$ de
$\Gamma(U_\lambda\cap U_\mu,\mathscr{G})$ tel que
\[
  \gamma_{\lambda\mu}\cdot
  (p_\mu|U_\lambda\cap U_\mu)=p_\lambda|U_\lambda\cap U_\mu;
\]
en outre, si $\lambda$, $\mu$, $\nu$ sont trois indices tels que
$U_\lambda\cap U_\mu\cap U_\nu\neq\varnothing$, les restrictions
$\gamma'_{\lambda\mu}$, $\gamma'_{\mu\nu}$, $\gamma'_{\lambda\nu}$ de
$\gamma_{\lambda\mu}$, $\gamma_{\mu\nu}$, $\gamma_{\lambda\nu}$ à
$U_\lambda\cap U_\mu\cap U_\nu$ vérifient la condition
$\gamma'_{\lambda\nu}=\gamma'_{\lambda\mu}\gamma'_{\mu\nu}$; autrement dit
$(\lambda,\mu)\mapsto\gamma_{\lambda\mu}$ est un \emph{1-cocycle} du
recouvrement $(U_\lambda)$ à valeurs dans $\mathscr{G}$. Si, pour tout
$\lambda$, $p'_\lambda$ est un second élément de
$\Gamma(U_\lambda,\mathscr{P})$, il existe un élément et un seul
$\beta_\lambda\in\Gamma(U_\lambda,\mathscr{G})$ tel que
$p'_\lambda=\beta_\lambda\cdot p_\lambda$ et le 1-cocycle
$(\gamma'_{\lambda\mu})$ correspondant à la famille $(p'_\lambda)$ est donné
par
\[
  \gamma'_{\lambda\mu}=\beta_\lambda\gamma_{\lambda\mu}\beta_\mu^{-1},
\]
autrement dit, est \emph{cohomologue} à $(\gamma_{\lambda\mu})$.
Inversement, la donnée d'un 1-cocycle $(\gamma_{\lambda\mu})$ définit, pour
tout couple $(\lambda,\mu)$, un automorphisme $\theta_{\lambda\mu}$ du
faisceau \emph{d'ensembles}
$\mathscr{G}|U_\lambda\cap U_\mu$, savoir la translation à droite par
$\gamma_{\lambda\mu}$, et le fait qu'il s'agisse d'un cocycle montre que l'on
peut \emph{recoller} les faisceaux d'ensembles $\mathscr{G}|U_\lambda$ à
l'aide des automorphismes $\theta_{\lambda\mu}$ (0, 3.3.1); on obtient ainsi
de façon évidente un torseur sous $\mathscr{G}$, soit $\mathscr{P}$, et si
l'on prend pour $p_\lambda$ la section unité au-dessus de $U_\lambda$, le
1-cocycle correspondant n'est autre que le 1-cocycle donné
$(\gamma_{\lambda\mu})$; en outre, si l'on remplace
$(\gamma_{\lambda\mu})$ par un 1-cocycle
$\gamma'_{\lambda\mu}=\beta_\lambda\gamma_{\lambda\mu}\beta_\mu^{-1}$ qui
lui est cohomologue, on vérifie aussitôt que le torseur obtenu est isomorphe
à $\mathscr{P}$.

En particulier, si $(\gamma_{\lambda\mu})$ est un \emph{1-cobord}, autrement
dit de la forme
$\gamma_{\lambda\mu}=\beta_\lambda\beta_\mu^{-1}$, le torseur obtenu
$\mathscr{P}$ est \emph{isomorphe à $\mathscr{G}$} (considéré comme torseur
sous lui-même pour les translations à gauche); on dit dans ce cas que
$\mathscr{P}$ est \emph{trivial}, et la réciproque est évidente.

Plus particulièrement, il résulte de (III, 1.3.1) que l'on a:
\end{env}

\begin{proposition}[16.5.16]
\label{IV.16.5.16}
Soient $Z$ un schéma affine, $\mathscr{G}$ un $\mathcal{O}_Z$-Module
quasi-cohérent; alors tout torseur sous $\mathscr{G}$ est trivial.
\end{proposition}

